%% V2 copy of the design rule, boxed in Background so Binding and Hiding can
%% instantiate it without forward reference.  This box states the rule, and the
%% justification follows it in Background.

\begin{rulebox}
\noindent\textbf{The design rule.}\quad
A protocol designer supplies five things: a field~$\field$, a security
level~$\secpar$, a hash~$H$, a choice of standard or hiding interface, and an
adversary model, declared as the query-budget exponent~$h$.  Everything else
is derived.  From $(\field,\secpar,h)$ the parameter constructor produces
\begin{equation}\label{eq:rule}
  D=\left\lceil\frac{\secpar+1+2h}{b}\right\rceil,\qquad
  k=\left\lceil\frac{\secpar}{b}\right\rceil,\qquad
  S=\left\lceil\frac{\secpar}{8}\right\rceil,
\end{equation}
where $b=\lfloor\log_2|\field|\rfloor$ is the conservative bits-per-element.
Here $D$ is the digest width in field elements, and $k$, $S$ are the field-
and byte-salt cardinalities.  Each output is a checkable obligation.  The
constructor either meets it and emits the derived parameters, or it refuses.
A designer can therefore neither quietly under-specify nor silently
over-claim.
\end{rulebox}
